\documentclass[./examen.tex]{subfiles}
\usepackage{csvsimple}
\usepackage[first=1, last=20]{lcg}

\newcommand{\random}{\rand\arabic{rand}}
\usepackage{verbatim}
\usepackage{moreenum}
\newcounter{z}
\newcommand\y[2]{
  \loop \ifnum\value{z} < #1
    #2%
    \stepcounter{z}%
  \repeat
}

\newcommand{\shuffle}[1]{%
    \foreach \i in {#1} {%
        \edef\temps{\noexpand\item \csname question\i\endcsname}%
        \temp
    }%
}

\begin{document}



\csvreader[
    head to column names, % Asigna nombres de columna
    separator=comma % Define la coma como el separador del CSV
    ]{./alumnos.csv}{}{% Itera sobre cada fila del CSV
 
\begin{center}

\makebox[\textwidth]{Esc:4-117 Ejército de los Andes, Evaluación Termodinámica. Nombre y Apellido:\nombre}
\end{center}

\begin{questions}
 \begin{randomize}     
     \question[1]  Determinar cuantos grados Kelvin representan \random3°F y cuántos °F representan $\random58 ^\circ C$

    \question[2]  Determinar cuantas calorías representan $4\random BTU$ y cuántos $Chu$ representan $\random58 BTU$
     
     
    \question[3] Calcular la cantidad de calor que suministra un bloque de \random0 Kg de material con calor específico medio $C_m=0,\random kcal/kg^\circ C$ cuya temperatura desciende de $\random0^\circ C$ a $\random^\circ C$.
    
     
    \question[4]   Calcular la cantidad de calor que se transmite por conductividad durante $\tau=1.5h$ a través de una losa de hormigón armado $C_t=1,3\random\frac{kcal m}{m^2 h ^\circ C}$ con una superficie de $S=\random 6m^2$ y un espesor $e=1\random cm$ cuando la temperatura exterior es de $t_e=\random05^\circ C$ y la interior es $t_i=\random^\circ C$
    \end{randomize}
 \end{questions}

}
\end{document}
